% Generated by tools/edn_latex.py from reports/edn.md.
% Do not edit: change reports/edn.md and run `make edn`.
\ednentry{
  id = {EDN-81},
  title = {Coding agents work under a written rule: an independent agent reviews, people decide and are accountable},
  date = {2026-09-29 (rule written down); review rule reworded on 2026-10-12, as a proposal},
  milestone = {Other (team process)},
  topic = {Use of AI Coding Agents, Code Review},
  participants = {@lukas2510 as Scrum master wrote the rule down and proposed the rewording, to be confirmed by the team at the laboratory session on 2026-10-13; no discussion of the rule by the other four members is on record.},
  decision = {Coding agents work under the four rules of the \href{https://github.com/mlops-2627q1-mds-upc/MLOPS_Recommenditos/blob/main/docs/docs/scrum/working-agreements.md}{working agreements}, and the review rule was reworded on 2026-10-12 to: ``Agent-written code goes through the same pull request as anything else. An independent agent with a fresh context reviews every substantive change by reproducing its claims and trying to break it, and the required checks gate the merge. A person decides the questions that belong in the EDN and is accountable for the change.''
It used to read ``The reviewer is a person, and reviewing means having understood the tests, not having watched them pass'', which the trace does not show.
No person signs off each pull request.},
  rationale = {\emph{Alternatives considered:}
\begin{itemize}[leftmargin=*, topsep=2pt]
\item \textbf{Option A (chosen): reword the review rule to the process as practised.}
\newline Pros: rule and practice agree. Cons: for most pull requests the reviewer is an agent and not a person.
\item \textbf{Option B: keep the old wording and record the gap as a weakness.}
\newline Pros: no rule changes before the team has met. Cons: a rule the team does not follow.
\item \textbf{Option C: a second person approves every pull request} (considered on 2026-09-21 and 2026-10-12).
\newline Pros: human review is visible on every pull request. Cons: GitHub blocks approving one's own pull request, and an agent's runs under its author's account (\href{https://github.com/mlops-2627q1-mds-upc/MLOPS_Recommenditos/pull/8}{PR \#8}), so most merges would wait.
\end{itemize}
\par
\emph{Why:} Of the 56 pull requests merged by 2026-10-11, 2 carry a review by another person, both by @kadameit, and branch protection requires 9 checks and 0 approvals.
The bodies of \href{https://github.com/mlops-2627q1-mds-upc/MLOPS_Recommenditos/pull/50}{\#50}, \href{https://github.com/mlops-2627q1-mds-upc/MLOPS_Recommenditos/pull/53}{\#53}, \href{https://github.com/mlops-2627q1-mds-upc/MLOPS_Recommenditos/pull/55}{\#55}, \href{https://github.com/mlops-2627q1-mds-upc/MLOPS_Recommenditos/pull/62}{\#62}, \href{https://github.com/mlops-2627q1-mds-upc/MLOPS_Recommenditos/pull/71}{\#71} and \href{https://github.com/mlops-2627q1-mds-upc/MLOPS_Recommenditos/pull/80}{\#80} describe an independent or adversarial review, which was an agent's; otherwise, before 2026-10-12, the review left its trace in the agent session and not on GitHub.
Since 2026-10-12 its result is posted on the pull request before the merge, first on \href{https://github.com/mlops-2627q1-mds-upc/MLOPS_Recommenditos/pull/101\#issuecomment-6124299563}{\#101}.
\par
People decide: of the 76 entries EDN-01 to EDN-78, 46 record the AI's proposal as accepted, 27 as modified, 2 as rejected and 1 as input for further analysis.
Comments and pull requests under a member's account are often typed by an agent on that member's instruction, and the member answers for them (the third rule).
The strongest cases:
@lukas2510 asked who owned the DagsHub repository the AI had recommended, which invalidated it (EDN-19);
he asked what availability the single unattended VM could actually meet, which replaced the draft requirements' 99\,\% uptime target (EDN-10, commit \texttt{3afcb50c} of \href{https://github.com/mlops-2627q1-mds-upc/MLOPS_Recommenditos/pull/19}{\#19});
he overruled the second of two AI passes on an omitted equipment list, on a ground neither pass had weighed (EDN-61);
@kadameit required two changes on \href{https://github.com/mlops-2627q1-mds-upc/MLOPS_Recommenditos/pull/26}{\#26}, and ran the suite on Windows 11 for \href{https://github.com/mlops-2627q1-mds-upc/MLOPS_Recommenditos/pull/85}{\#85}, where his approval came three minutes after the merge and did not inform it;
@lukas2510 moved \href{https://github.com/mlops-2627q1-mds-upc/MLOPS_Recommenditos/pull/86}{\#86} back to draft, because accepting SC-04's miss before tuning was premature.
\par
Two examples of where checking failed: the DagsHub recommendation was posted on \#26 unchecked, and two pull requests moved the laboratory to the wrong weekday (corrected in \href{https://github.com/mlops-2627q1-mds-upc/MLOPS_Recommenditos/pull/97}{PR \#97}).
The weak spot: 48 of the 56 were opened and merged under one account with no second person's review, so in sprint 3 every teammate opens the pull request for their own ticket.},
  ai = {Information seeking, Alternative generation, Alternative assessment, Recommendation, Solution generation},
  response = {Accepted},
  assessment = {The oracle criterion and the rule's text are AI's, adopted by @lukas2510 on 2026-09-29.
On 2026-10-12 AI measured the numbers above, found that the review rule and the trace disagreed, and recommended option A over B, which @lukas2510 chose.
He decided against option C the same day; option C is AI's as well.},
  aievidence = {Claude Code sessions of 2026-09-21, 2026-09-29 (sprint 2 planning) and 2026-10-12.},
  evidence = {\href{https://github.com/mlops-2627q1-mds-upc/MLOPS_Recommenditos/pull/45}{PR \#45}, \href{https://github.com/mlops-2627q1-mds-upc/MLOPS_Recommenditos/issues/95}{issue \#95}; \texttt{gh pr list -\allowbreak{}-\allowbreak{}state merged -\allowbreak{}-\allowbreak{}json number,author,mergedBy,reviews} and the branch protection of \texttt{main}, both read on 2026-10-12.},
}
